\paragraph*{04.04.05.01. Initiative ergreifen und proaktiv mit Rat und Tat zur Seite stehen}
\kcilabel{para:04.04.05.01_InitiativeErgreifen}{04.04.05.01}
\label{para:04.04.05.01_InitiativeErgreifen}

%%%% Task
\par Als \gls{wsl} verstand ich meine Aufgabe als proaktive Führung in komplexem, heterogenem Projektumfeld. Angesichts hoher Schnittstellen-Abhängigkeiten, dynamischer Programmplanung und Parallelbetrieb-Herausforderungen (\gls{r3}/\gls{s4}) war frühzeitige Initiative entscheidend für Risikobeherrschung, Systemstabilität und Koordination.

%%% Action
\par Entwickelte Konzept für Testdaten-Beschaffung aus Produktivbetrieb und Synchronisationsplan für Parallelbetrieb. Nutzte proaktives Risiko-/Eskalationsmanagement für Datenverarbeitung in \gls{s4}. Setzte intensive Kommunikation ein für Informationsaustausch und Akzeptanzförderung (vgl. \autoref{fig:R2HStatusSSTJuni2025}).

%Result
\par Maßnahmen konnten nicht alle Herausforderungen beseitigen; Führung musste kontinuierlich angepasst werden. Enge Zusammenarbeit mit Partnern und Stakeholdern sicherte Stabilität und Zielerreichung. Risiken in \gls{s4}-Datenverarbeitung führten zu Verzögerungen.

%Reflect
\par Zukünftig: Bündelung in einem Projekt und frühzeitige Anforderungen (Kontextdiagramm, Schnittstellenbeschreibung) reduzieren Parallelbetrieb-Risiken. Proaktive Führung war entscheidend für Ziele und Stabilität.